\begin{cnabstract}

在社交媒体中，信息传播受内容特征、主体行为与网络结构的共同驱动。近年来，大语言模型驱动的AI代理已在多个社交媒体平台独立参与讨论，与人类用户共存于同一传播系统，形成多主体共存的传播格局。然而，现有传播分析与级联预测方法普遍假设传播节点同质，未将主体类型纳入建模；公开数据集亦缺乏主体标注，制约了多主体传播研究。针对上述问题，本文设计并实现了多主体信息传播影响力评估系统，从消息层、路径层和交互类型层三个维度量化评估人类与AI代理的传播影响力差异。本文所研究的传播影响力指信息内容在传播网络中引发后续回应与持续扩散的能力，涵盖内容吸引力与扩散持续力两个方面。主要工作如下：

（1）提出基于组感知融合的评论传播价值评估模型。该模型将主体标识作为显式特征组引入，通过门控融合机制自适应调整发布环境、文本语义、情感倾向和主体标识四组异构特征的贡献度，实现评论级传播价值评估。

（2）提出基于交互感知的异构路径传播延续预测模型。该模型以评论树中从根节点到目标节点的路径为输入，通过交互类型编码层显式建模四种主体交互模式，结合主体感知节点编码与LSTM序列编码，实现路径级传播延续预测。

（3）设计并实现完整的评估系统，采用四层架构与六个功能模块，涵盖数据导入与预处理、模型推理、多主体对比分析及可视化展示，将两个模型的输出按主体类型聚合，量化两类主体的传播影响力差异。系统验证结果表明，系统主要业务流程可正常运行，核心功能符合设计预期。

    \cnkeywords{多主体信息传播；传播影响力评估；深度表格学习；异构路径编码}
\end{cnabstract}
